\subsection{Limites et critiques}

\begin{itemize}
	
	\item Limites générales de l'IA
	
les modèes d'IA sont entraînés en anglais, pose des questions d'adaptabilité : "Les IA conversationnelles reposent sur des grands modèles de langue (LLM) entraînés principalement sur des données en anglais, ce qui crée des biais linguistiques et culturels dans les résultats qu'ils produisent. [...] Il est donc indispensable que ces modèles soient entrainés avec une diversité de sources culturelles et
linguistiques, y compris les langues régionales et les expressions artistiques moins représentées" \footcite{2025a}
	
	\item Limites contextuelles 

«How to give context to the machines? To give common sense? It's a problem of artificial intelligence from the beginning.»
\footcite{2021}

\item limites de l'automatisation :
 
"l’intégration de la vision artificielle aux méthodes des historiens n’existe que conjointement à des
interventions humaines, qui assurent la pertinence des traitements effectués et apportentune analyse nécessaire." \footcite{norindr2023}

\item sur l'indexation automatique audio :
« Audio archives present a complex terrain for contemporary speech processing technologies, owing to the varied domains these recordings encapsulate. » \footcite{sullivan2025} : deux problèmes sur un corpus d'archives anciennes : variabilité linguistiques (multilingues, accents différents), variabilité temporelle (vieillissement des voix).
les archives audiovisuelles sont un matériau hétérogène : opacité du médium. 

\item sur l'indexation automatique des archives audiovisuelles 

IL est important pour les archivistes de savoir comment ont été entraînés les modèles qu'ils utilisent, sur quelles données.

la vidéo : souvent de moins bonne qualité que l'image fixe, donc les résultats d'indexation automatique peuvent être moins bons

contrairement aux documents textuels qui comportent souvent des éléments d'indexation dans leur forme (couverture de livre, ) les archives audiovisuelles n'en contiennent peu ou pas, ce qui les rend opaques et peu réceptives aux pratiques d'indexation. Il s'agit alors de mettre en place une « indexation par le contenu ». \footcite{zotero-item-582} 

\subitem Lacunes 

l'article de attard souligne que  « The study underscores the lack of specialised annotation tools and the high resource demands of sophisticated models, which could present significant obstacles for future research with limited computational resources. » \footcite{attard2025} c'est une limite surtout pour les isntitutions. Rejoint l'argument et les limites exposées par \footcite{sullivan2025}
 
Les modèles multimodaux qui combinent audio et vidéo (vivit-ast) ne donnent pas les meilleurs résultats. Il faudrait donc adopter des architectures plus simples, qui demandent moins de calcul d'usage et qui séparent les pipelines. \footcite{attard2025}

Absence d'outils d'annotation spécialisés dans la segmentation vidéo. Attard adapte label studio, un outil généraliste.
 \footcite{attard2025}
 
Limites des transformers :
« ViT has been shown to only be effective when trained on large-scale datasets, as transformers lack some of the inductive biases of convolutional networks, and thus require datasets larger than the common ImageNet ILSRVC dataset to train. » \footcite{arnab2021}. : les transformers n'ont pas les biais des CNN qui leur permettraient de faire des généralisations et de reconnaître un objet où qu'il soit dans l'image. Ils ont donc besoin d'être entrainés sur beaucoup plus de données : c'est une grande limite. Une approhe par calssifieurs d'images plus légers comme ResNet peut donc être préférable. ViViT obtient de très bons résultats mais entraînés sur des corpus vidéos contemporains : vidéos youtube. Diffiilement adaptable à un corpus d'archives historique sauf pré entraînement massif. 


\item éthique de l'IA

Sur l'éthique de l'exprience utilisateur : Principles for Accountable Algorithms : responsability, explainability, accuracy, auditability, fairness
\footcite{zotero-item-574}

	Limites du machine learning : 
"Une répartition arbitraire des systèmes est toutefois limitée dans la mesure où, d’une part, les systèmes entraînés ont souvent un domaine d'application spécifique
et, d’autre part, leurs capacités sont fortement liées aux jeux de données utilisés. Les algorithmes entraînés avec différents jeux de données produisent des résultats
différents."
\footcite{zotero-item-584}

\item Biais de la computer vision appliquée aux images : 

"For instance, racial bias has been an ongoing challenge with	computer vision tools, particularly with facial recognition software. A study	published by MIT Media Lab in 2018 found that facial recognition software	had an average error rate of 0.8 percent for light-skinned men, versus a 34.7 percent	error rate for darker-skinned women" \footcite{fewster}
"societal biases are often engrained in	AI model training datasets, which lead to the models reproducing those	biases in their outputs" \footcite{fewster}
"In addition to this, the training data of computer	vision models has relied mostly on visual content available to be harvested	from the Internet, which in its majority corresponds to content developed	during the last three decades. These models would underperform on images	that are of a more historical character, commonly found in visual archives,	since they are underrepresented in the training data" \footcite{fewster}

Limite liée à la qualité des images : 
la qualité des images est absolument primordiale et déterminante dans la mise en place d'un pipeline d'indexation automatique : "La question de la qualité des images est
cruciale dans le cadre d’un projet de recherche en vision artificielle, puisque celle-ci a la possiblité d’impacter les performances de l’algorithme utilisé, et d’affecter négativement les interprétations du modèle". \footcite{norindr2023}
Important dans le cadre d'une indexation d'archives audiovisuelles, qui sotn souvent de moins bonne qulité que les images car les fichiers sont plus lourds.

limites des résraux de neurones profonds tels que les CNN pour l'analyse de l'esthétisme des images : « While hand-crafted features allow to quantify which feature contributes the most to an aesthetic rating, this interpretability is lost with generic features. [...] Deep neural networks and generic features basically represent a black-box approach that lacks this kind of interpretability. » \footcite{brachmann2017}
"with deep learning models, their success comes at a cost. At present, the understanding of deep features and how they work in object or aesthetic recognition lacks behind. » \footcite{brachmann2017}
Le problème est le suivant avec les CNN, les décisions qui sont prises sont obscures. Pour l'indexation des archives, c'est problématique si les archivistes ne sont pas capables d'expliquer pourquoi telle classification a été choisie par le modèle. Il faut maintenir une supervision humain dans le pipeline d'indexation. 
La question de l'interopérabilité et de la transparence est aussi importante que la question de performance en sciences humaines et dans le milieu du patrimoine, car il s'agit de donner accès aux connaissances. (à relier avec l'éthique de l'IA : principes MA/ML : explainability)

La machine analyse différemment de l'humain, de manière objective et précise et exhaustive, qu'apporte l'humain ? : "the computational analysis of paintings has several advantages compared to an analysis carried out by human experts. For example, a computational analysis can pick up very subtle relationships that may escape the attention by human observers; moreover, computational methods are objective in nature and are potentially non-exhaustive in the amount of detail analyzed. »\footcite{brachmann2017}

\item Dépendance à la qualité de l'image

Limite exprimée dans le cadre de l'utilisation de line segment detection : 
"The first general challenge for line segment detection and description is that the endpoints of detected line segments are generally unstable due to many factors, such as image noise, occlusion, and illumination variance. This makes the line segment description and matching difficult since the corresponding line ROS (Region of Support) or neighborhood relationship among line segments is generated according to the line segment endpoints. »\footcite{lin2024}
La compréhension et l'analyse de l'image dépend de la qualité de l'image, souvent réduite pour des raison de stockage de données, pendant l'entraînement. 

\item FRBR
le FRBR mentionne les images animées
mais pas le son en tant que composante principale


est ce qu'il u a des choses qu'on ne peut pas automatiser ? que les fichiers numériques dont l'indexation peut être automatisée. or, les colelctiosn actuelles comportent bcp de bobines : les lire suppose la dégradation du support. 

sur l'IA chez skinsoft : 
le moteur doit il être entraîné sur des donnes publiques, ou entraîné complémentairement par des données de l'institution sur laquelle il travaille. il faut que les clients acceptent : contrat spécifique. (limites) 

clients qui ne veulent pas de l'IA : prévoir une application qui tourne toujours sans IA : comment faire ça ? la proposer que pour certains clients 

Il faut l'accord client et les données d'un seul client ne sont pas suffisantes.

\item la tension privé public
Si skinsoft est une entreprise privée, elle adopte en revanche une posture publique culturelle, celle de ses clients, afin de répndre à leurs besoins. La mise en place de l'IA est pensée pour améliorer les logiques patrimoniales quotidiennes telles que l'indexation, le catalogage, le classement, la gestion de projet... Ces pratiques ne sont pas liées à des objectifs de rentabilité et de productivité concurrentiels comme c'est le cas d'entreprises privées. Ainsi, skinsoft fait face à une tension : l'entreprise s'affiche comme une entreprise innovante en matière de technologies numériques, alors même que ses clients sont cantonnés à des pratiques presqu'ancestrales, celle de la gestion de collections. 

Au sein de la mise en place et de l'arrivée de l'IA dans un logiciel de gestion des collections, il faut ainsi se questionner sur sa pertinence. Ces nouvelles fontionnalités sont elles attendues ? Les institutions sont-elles prêtes à utiliser l'IA au quotidien pour gérer leur collection? pour les décrire ? pour les rendre accessible au public ?

En 2025, le ministère de la culture identifie "l'enrichissement de la connaissances et des savoirs" comme un usage potentiel de l'IA en secteur culturel avec par exemple "de l’aide au catalogage, à l’indexation automatique des collections, au tri et au classement des
documents et données destinés à être archivés, à la transcription des textes manuscrits, à la reconnaissance de formes, à la diffusion et à la valorisation des contenus culturels et patrimoniaux quel que soit leur
format (image, son…)". \footcite{2025a}

Les limites de l'IA identifiées, à savoir l'hallucination, la marge d'erreur, les biais culturels, sont ici restreintes puisque les outils d'IA seraient mis à disposition apr skinsoft à des professionnels, et non au grand public directement. Certes, ces outils IA permettront à terme, comme l'indexation automatique, de rendre les collectiosn acessibles au public, mais seulement après vérification d'un professionnel attitré. 

Les professionnels pourraient mettre à profit l'IA pour un double objectif : la préservation et l'accessibilité des collections. \footcite{zotero-item-640}. C'est exactement le cas pour l'indexation automatique. 

limites indexation : 
Moiraghi et Diginpix : "les bases de données d'indexation constituées à partir de la reconnaissance des entités sont conséquentes (elles peuvent avoir une taille supérieure à celle des vidéos elles-mêmes)". cela suppose que l'utilisateur doivent vérifier un nombre de données assez important. 

\end{itemize}